\begin{remark}
\label{remark-map-into-derived-limit-truncations}
Let $\mathcal{A}$ be an abelian category. Let $K^\bullet$ be a
complex of $\mathcal{A}$. Then $\tau_{\geq -n}K^\bullet$ is an
inverse system of complexes and which in particular determines
an inverse system in $D(\mathcal{A})$. Let us assume that
$R\lim \tau_{\geq -n}K^\bullet$ exists. Then the canonical maps
$c_n : K^\bullet \to \tau_{\geq -n}K^\bullet$ are compatible with the
transition maps of our inverse system. By the defining distinguished
triangle of Definition \ref{definition-derived-limit} and
Lemma \ref{lemma-representable-homological}
we conclude there exists a morphism
$$
c : K^\bullet \longrightarrow R\lim \tau_{\geq -n}K^\bullet
$$
in $D(\mathcal{A})$ such that the composition of $c$ with the projection
$R\lim \tau_{\geq -n}K^\bullet \to \tau_{\geq -m}K^\bullet$ is equal to $c_m$.
Now the morphism $c$ may not be unique, but we claim that whether or not
$c$ is an isomorphism is independent of the choice of $c$ (and of our choice
of the homotopy limit). Namely, for $i \in \mathbf{Z}$ and for $m > -i$
the composition
$$
H^i(K^\bullet) \xrightarrow{H^i(c)} H^i(R\lim \tau_{\geq -n}K^\bullet)
\to H^i(\tau_{\geq -m}K^\bullet) = H^i(K^\bullet)
$$
is the identity. Hence $H^i(c)$ is an isomorphism if and only if the
second map is an isomorphism. This is independent of $c$ and also
independent of the choice of the homotopy limit (as any two choices
are isomorphic).
\end{remark}